\section*{Capítulo \ref{ch:g10:the-mole} --- O mol e a massa molar}

\begin{solution}{exo:g10:the-mole:1}
$M(\ce{O2}) = 2 \times 16.0 = \qty{32.0}{g/mol}$;
$M(\ce{N2}) = 2 \times 14.0 = \qty{28.0}{g/mol}$;
$M(\ce{CH4}) = 12.0 + 4 \times 1.0 = \qty{16.0}{g/mol}$;
$M(\ce{NH3}) = 14.0 + 3 \times 1.0 = \qty{17.0}{g/mol}$;
$M(\ce{CO2}) = 12.0 + 2 \times 16.0 = \qty{44.0}{g/mol}$.
\end{solution}

\begin{solution}{exo:g10:the-mole:2}
Água: $n = 36.0 / 18.0 = \qty{2.00}{mol}$. Metano:
$n = 4.0 / 16.0 = \qty{0.25}{mol}$.
\end{solution}

\begin{solution}{exo:g10:the-mole:3}
$m(\ce{NaCl}) = 0.25 \times 58.5 = \qty{14.6}{g}$, cerca de \qty{15}{g};
$m(\ce{CO2}) = 2.0 \times 44.0 = \qty{88}{g}$.
\end{solution}

\begin{solution}{exo:g10:the-mole:4}
$N = 0.50 \times \num{6.02e23} = \num{3.0e23}$ \omterm{def:g7:atoms-and-molecules:molecule}{moléculas} de gás oxigênio.
$n = \num{3.0e24} / \qty{6.02e23}{mol^{-1}} = \qty{5.0}{mol}$ de ferro.
\end{solution}

\begin{solution}{exo:g10:the-mole:5}
$V = n \times V_m = 0.50 \times 24.1 = \qty{12}{L}$.
$n = V / V_m = 12 / 24.1 = \qty{0.50}{mol}$.
\end{solution}

\begin{solution}{exo:g10:the-mole:6}
A gota tem massa de \qty{0.050}{g}, logo
$n = 0.050 / 18.0 = \qty{2.8e-3}{mol}$ e
$N = \num{2.8e-3} \times \num{6.02e23} = \num{1.7e21}$ \omterm{def:g7:atoms-and-molecules:molecule}{moléculas}.
\end{solution}

\begin{solution}{exo:g10:the-mole:7}
Um \omterm{def:g7:atoms-and-molecules:atom}{átomo} de alumínio (\qty{27.0}{g/mol}) é mais ou menos metade tão
pesado quanto um \omterm{def:g7:atoms-and-molecules:atom}{átomo} de ferro (\qty{55.8}{g/mol}), por isso a mesma
massa de alumínio contém cerca de duas vezes mais \omterm{def:g7:atoms-and-molecules:atom}{átomos}. Verificação:
alumínio $10/27.0 = \qty{0.370}{mol}$, isto é,
\num{2.23e23} \omterm{def:g7:atoms-and-molecules:atom}{átomos}; ferro $10/55.8 = \qty{0.179}{mol}$, isto é,
\num{1.08e23} \omterm{def:g7:atoms-and-molecules:atom}{átomos}. O alumínio ganha, por um fator de cerca de 2.1.
\end{solution}

\begin{solution}{exo:g10:the-mole:8}
Um litro é $1/24.1 = \qty{0.0415}{mol}$ de gás. Dióxido de carbono:
$0.0415 \times 44.0 = \qty{1.83}{g}$; gás nitrogênio:
$0.0415 \times 28.0 = \qty{1.16}{g}$. O dióxido de carbono é cerca de
uma vez e meia mais denso que o \omterm{def:g4:air-a-mixture-of-gases:air}{ar}, cuja \omterm{def:g10:the-mole:molar-mass}{massa molar} é próxima da do gás
nitrogênio. O suco em fermentação libera dióxido de carbono, que desce e
se acumula perto do chão de uma adega fechada, onde pode asfixiar quem
entrar.
\end{solution}

\begin{solution}{exo:g10:the-mole:9}
$M(\ce{C6H12O6}) = 6 \times 12.0 + 12 \times 1.0 + 6 \times 16.0 =
\qty{180.0}{g/mol}$. Seta $\div N_A$: $n = \num{1.0e22} / \num{6.02e23} =
\qty{1.66e-2}{mol}$; seta $\times M$: $m = \num{1.66e-2} \times 180.0 =
\qty{3.0}{g}$.
\end{solution}

\begin{solution}{exo:g10:the-mole:10}
$M = \qty{342.0}{g/mol}$, logo $n = 5.0 / 342.0 = \qty{1.46e-2}{mol}$;
$N = \num{1.46e-2} \times \num{6.02e23} = \num{8.8e21}$ \omterm{def:g7:atoms-and-molecules:molecule}{moléculas}. Cada
uma tem 12 \omterm{def:g7:atoms-and-molecules:atom}{átomos} de carbono: $12 \times \num{8.8e21} = \num{1.1e23}$ \omterm{def:g7:atoms-and-molecules:atom}{átomos} de carbono.
\end{solution}

\begin{solution}{exo:g10:the-mole:11}
Água: $1000 / 18.0 = \qty{55.6}{mol}$. Etanol:
$M = 2 \times 12.0 + 6 \times 1.0 + 16.0 = \qty{46.0}{g/mol}$, logo
$1000 / 46.0 = \qty{21.7}{mol}$. A água tem a maior quantidade, por um
fator $55.6 / 21.7 \approx 2.6$, que é simplesmente $46.0 / 18.0$.
\end{solution}

\begin{solution}{exo:g10:the-mole:12}
Ouro: $\qty{3.75}{g}$, $n = 3.75 / 197.0 = \qty{1.90e-2}{mol}$, isto é,
\num{1.15e22} \omterm{def:g7:atoms-and-molecules:atom}{átomos}. Cobre: \qty{1.25}{g},
$n = 1.25 / 63.5 = \qty{1.97e-2}{mol}$, isto é, \num{1.19e22} \omterm{def:g7:atoms-and-molecules:atom}{átomos}.
Surpreendentemente, o anel contém um pouco mais \omterm{def:g7:atoms-and-molecules:atom}{átomos} de cobre que de
ouro: um \omterm{def:g7:atoms-and-molecules:atom}{átomo} de ouro é cerca de três vezes mais pesado que um \omterm{def:g7:atoms-and-molecules:atom}{átomo} de
cobre.
\end{solution}

\begin{solution}{exo:g10:the-mole:13}
\begin{enumerate}
  \item $V = 8.0 \times 6.0 \times 3.0 = \qty{144}{m^3} =
        \qty{1.44e5}{L}$; $n = \num{1.44e5} / 24.1 = \qty{5.98e3}{mol}$;
        $N = \num{5.98e3} \times \num{6.02e23} = \num{3.6e27}$ \omterm{def:g7:atoms-and-molecules:molecule}{moléculas}.
  \item $M = (78 \times 28.0 + 21 \times 32.0 + 1 \times 40.0) / 100 =
        2896 / 100 = \qty{29.0}{g/mol}$.
  \item $m = \num{5.98e3} \times 29.0 = \qty{1.73e5}{g}$, cerca de
        \qty{173}{kg}: o \omterm{def:g4:air-a-mixture-of-gases:air}{ar} de uma sala de aula pesa tanto quanto dois
        ou três adultos.
\end{enumerate}
\end{solution}

\begin{solution}{exo:g10:the-mole:14}
\begin{enumerate}
  \item $n = 1000 / 28.1 = \qty{35.6}{mol}$, logo
        $N = 35.6 \times \num{6.02e23} = \num{2.14e25}$ \omterm{def:g7:atoms-and-molecules:atom}{átomos}.
  \item $m = \qty{1.000}{kg} / \num{2.14e25} = \qty{4.67e-26}{kg}$.
  \item A massa da esfera e a \omterm{def:g10:the-mole:molar-mass}{massa molar} dão a sua \omterm{def:g10:the-mole:amount}{quantidade de
        matéria}, $n = m / M$. Se os \omterm{def:g7:atoms-and-molecules:atom}{átomos} $N$ forem contados de modo
        independente (a partir do volume da esfera e do espaçamento dos
        \omterm{def:g7:atoms-and-molecules:atom}{átomos} no cristal), a razão $N / n$ é a \omterm{def:g10:the-mole:avogadro}{constante de Avogadro}.
\end{enumerate}
\end{solution}

\begin{solution}{exo:g10:the-mole:15}
$M = 0.758 \times 35.0 + 0.242 \times 37.0 = 26.53 + 8.95 =
\qty{35.5}{g/mol}$: o valor da tabela.
\end{solution}

\begin{solution}{pb:g10:the-mole:1}
\textbf{1.} $M(\ce{H2O}) = 2 \times 1.0 + 16.0 = \qty{18.0}{g/mol}$.

\textbf{2.} $m = 250 \times \qty{1.0}{g} = \qty{250}{g}$.

\textbf{3.} $n = 250 / 18.0 = \qty{13.9}{mol}$.

\textbf{4.} Cada \omterm{def:g7:atoms-and-molecules:molecule}{molécula} tem dois \omterm{def:g7:atoms-and-molecules:atom}{átomos} de hidrogênio e um \omterm{def:g7:atoms-and-molecules:atom}{átomo} de
oxigênio: \qty{27.8}{mol} de \omterm{def:g7:atoms-and-molecules:atom}{átomos} de hidrogênio, \qty{13.9}{mol} de
\omterm{def:g7:atoms-and-molecules:atom}{átomos} de oxigênio.

\textbf{5.} $N = 13.9 \times \num{6.02e23} = \num{8.36e24}$ \omterm{def:g7:atoms-and-molecules:molecule}{moléculas}.

\textbf{6.} $m = \qty{18.0}{g/mol} / \qty{6.02e23}{mol^{-1}} =
\qty{2.99e-23}{g}$.

\textbf{7.} $\num{8.36e24} \times \qty{2.99e-23}{g} = \qty{250}{g}$, a
massa da água do copo, como deveria ser.

\textbf{8.} $\num{8.36e24} / \num{3.16e7} = \num{2.6e17}$ anos: um
número de anos com dezessete zeros, muito além de qualquer contagem
possível.

\textbf{9.} $\qty{1.335e9}{km^3} \times \num{e9} = \qty{1.335e18}{m^3}$,
e com $\qty{1}{m^3} = \qty{1000}{L}$, \qty{1.335e21}{L}.

\textbf{10.} $\num{8.36e24} / \num{1.335e21} = \num{6.26e3}$ \omterm{def:g7:atoms-and-molecules:molecule}{moléculas}
marcadas por litro.

\textbf{11.} $\num{6.26e3} \times 0.250 = \num{1.57e3}$ \omterm{def:g7:atoms-and-molecules:molecule}{moléculas}
marcadas no segundo copo.

\textbf{12.} Massa $\num{1.335e21} \times \qty{1.0}{kg} =
\qty{1.335e21}{kg} = \qty{1.335e24}{g}$; $n = \num{1.335e24} / 18.0 =
\qty{7.42e22}{mol}$.

\textbf{13.} $13.9 / \num{7.42e22} = \num{1.87e-22}$.

\textbf{14.} $\num{1.87e-22} \times \num{8.36e24} = \num{1.57e3}$: a
mesma resposta da questão 11, como tem de ser.

\textbf{15.} $\num{1.335e21} / 0.250 = \num{5.34e21}$ copos.

\textbf{16.} $\num{8.36e24} / \num{5.34e21} \approx 1570$. Um copo
contém cerca de 1600 vezes mais \omterm{def:g7:atoms-and-molecules:molecule}{moléculas} do que os oceanos contêm copos
de água; espalhadas por igual, as \omterm{def:g7:atoms-and-molecules:molecule}{moléculas} de um copo deixam cerca de
1600 em cada copo de oceano.

\textbf{17.} $1 / \num{6.26e3} = \qty{1.6e-4}{L} = \qty{0.16}{mL}$,
cerca de três gotas.

\textbf{18.} O volume dos oceanos é uma estimativa conhecida, no máximo,
com três ou quatro algarismos, e foi arredondado; a água do mar não é
água pura (ela contém sais); o copo não tem exatamente \qty{250}{mL}; e
a \omterm{def:g6:pure-substances-and-mixtures:mixture}{mistura} por todos os oceanos nunca seria perfeitamente uniforme. Só a
ordem de grandeza merece confiança.

\textbf{19.} Cerca de 1600 \omterm{def:g7:atoms-and-molecules:molecule}{moléculas} do primeiro copo voltam no segundo.
\end{solution}
